\section{Discussion}

\subsection{Limitations}

\textbf{Persona as approximation}. Our personas are simplified models of real researchers. They capture research focus and characteristic questions but cannot replicate the full depth of domain knowledge, personal experience, or evolving research interests.

\textbf{Model dependence}. Calibration results may vary across LLM providers and versions. Our analysis uses a single model family; cross-model validation is needed.

\textbf{Ethical considerations}. Using real researcher identities for persona construction raises questions about consent and representation. We use publicly available information (publications, affiliations) and do not claim to represent these researchers' actual opinions.

\subsection{Distinctness vs.\ Quality}

A key concern is whether distinct reviews are also \emph{useful} reviews. We assess review quality on a sample of 40 reviews (8 per persona category) using three criteria: (1) Are major issues valid? (i.e., correctly identify a real weakness), (2) Are issues actionable? (provide specific guidance for improvement), (3) Are issues paper-specific? (reference specific content, not generic concerns).

\begin{table}[h]
\centering
\small
\begin{tabular}{lccc}
\toprule
Distinctness Tier & Valid (\%) & Actionable (\%) & Paper-Specific (\%) \\
\midrule
High ($|\rho| < 0.30$) & 78 & 72 & 81 \\
Medium ($0.30 \leq |\rho| \leq 0.60$) & 82 & 76 & 74 \\
Low ($|\rho| > 0.60$) & 79 & 71 & 65 \\
\bottomrule
\end{tabular}
\caption{Review quality by distinctness tier.}
\end{table}

\noindent Distinct reviewers produce reviews that are equally valid and actionable as convergent reviewers, while being significantly more paper-specific ($81\%$ vs.\ $65\%$). This confirms that distinctness does not come at the cost of quality.

\subsection{Practical Guidelines}

Based on our analysis:
\begin{itemize}
    \item \textbf{Always include key questions}: This is the single most effective calibration mechanism (+0.15 JS divergence)
    \item \textbf{Diversify across categories}: Select reviewers from at least 3 different expertise categories
    \item \textbf{Monitor convergence}: If all reviewers agree on everything, the panel lacks diversity
    \item \textbf{Calibrate to venue}: Include reviewers who would plausibly serve at the target venue
    \item \textbf{Match expertise to paper}: Personas reviewing outside their domain revert to generic feedback (Section 4.6)
    \item \textbf{Specify key questions concretely}: Specific questions (``What's the compute/memory tradeoff?'') outperform general ones (``Is this novel?'')
\end{itemize}

\subsection{Reproducibility}

All reviews were generated using Claude Sonnet 4.5 (\texttt{claude-sonnet-4-5-20250929}) with temperature 1.0 (default), no top-$k$ or top-$p$ overrides. The complete prompt template is:

\begin{lstlisting}
You are {name} from {affiliation}, an expert in {expertise}.
Your characteristic concern when reviewing papers is:
"{key_question}"

You are reviewing for {venue}, where papers are expected to
demonstrate {venue_standards}.

Review the following paper using the structured format below.
Score on a 1-4 scale where:
  1 = Reject, 2 = Weak Accept, 3 = Accept, 4 = Strong Accept

Required sections: Overall Assessment, Score, Major Issues,
Minor Issues, Strengths, Questions, Recommendations, Verdict.

[Full paper text follows]
\end{lstlisting}

\noindent Each review is generated independently --- reviewer $i$'s output is not visible to reviewer $j$. Average token count per review: approximately 1,200 tokens.
